\section{Locally certified state-fold search}
\label{sec:local-fold-search}

We implemented a bounded superoptimizer for literal binary Turing-machine
tables.  Its purpose was to find a further state reduction with a finite,
checkable simulation argument.  It found no reduction in the three search
classes described below, on either the published rotation240 table or the
new, repaired boundary240 table.  As a positive control, exactly the same
implementation rediscovers the established $90\to72$ reduction of the
rotated241 table.  These negative results concern the stated certificate
classes; they are not lower bounds on the number of states needed for the
consistency problem.

\subsection{A sound local invariant}

Fix a radius $K$.  At control state $q$, the abstraction retains the bits
in tape cells at offsets $-K,\ldots,K$ from the head.  We represent the
window by an integer $m$, with bit $i+K$ representing offset $i$.  The
abstract transition writes the center bit, shifts the window according
to the move direction, and nondeterministically assigns either value to
the newly exposed boundary cell.  Breadth-first closure computes a set
$I_K(q)$ of allowed masks for every state.

The initial set contains \emph{every} mask at the designated start state.
Consequently this invariant is valid from every initial binary tape, and
in particular from the blank tape.  This is a deliberate overapproximation:
the search does not assume that the source program, its register encoding,
or the startup routine supplies additional structure.  The abstraction
is sound because the actual newly exposed cell always has one of the two
values included in the abstract successor relation.  Induction on the
number of machine steps proves that every reachable concrete
configuration projects into $I_K$.

Soundness of this invariant must not be confused with completeness for
the application.  A mask in $I_K(q)$ need not arise from a well-formed
register frame, or from any execution of the fixed source program.
Such masks can rule out a proposed rewrite which would be valid under a
stronger invariant.  Enlarging the radius can sometimes recover a lost
correlation, but the results below are confined to $K=9$.

\subsection{Positive-block replacement domains}

For each row $(q,b)$, the search considers every possible replacement
edge $(w,d,t)$: either write bit, either move direction, and any existing
working state or the designated halt as target.  There are
$4(n+1)$ candidate edges per row for an $n$-state table.  Each candidate
is executed separately on every mask in $I_K(q)$ whose center is $b$.
After its first transition, execution follows the original table.

The candidate must reach a specified original configuration in at
least one and at most $F$ transitions.  The comparison includes the
control state, head displacement, and \emph{every bit of the initial
window}.  Neither trajectory may read a cell outside that window.
The final head position may be one cell beyond it when the final move
does not read that cell.  Since all changed cells lie inside the
window, equality of the compared windows proves equality of the entire
tapes, irrespective of the unrepresented cells.  Thus an accepted
domain is a universally quantified local fact, not a test of one
finite tape.

The implementation records which control states the replacement
executes strictly between its start and matching endpoint.  These
support sets prevent composition from silently using a row that the
proposed fold also changes.  All simulated old and new blocks have
positive length.  Matching checkpoints therefore preserve halting in
both directions and cannot introduce an infinite sequence of
zero-length simulation steps.

The three modes are as follows.

\begin{enumerate}
\item \textbf{Pair mode.}  The matching target is the configuration after
one original transition.  Intermediate replacement execution may not
execute state $q$.  For two states $q,r$, the algorithm intersects their
replacement-edge domains separately for read bits $0$ and $1$, then
requires both support sets to avoid both $q$ and $r$.  A nonempty
intersection for each read bit supplies common rows for the quotient
state.  Incoming targets and the designated start, if needed, are
mapped to that state.  The comparison uses the same tape and head at
each checkpoint; it does not permit an arbitrary tape representation
change or a permanently shifted head.

\item \textbf{Block mode.}  The original trajectory may take any positive
number of steps up to $F$ before the matching checkpoint, instead of
exactly one.  The replacement remains bounded by $F$ and uses the same
window.  The same support exclusions and pairwise domain intersection
are applied.  This allows a longer old block to be implemented by a
different positive block, but still requires an exactly matching
control state, head, and tape at the checkpoint.

\item \textbf{Incoming-deletion mode.}  For a state $s$ to be deleted,
every external row targeting $s$ must receive a replacement edge.
Each edge must rejoin a positive original block at a state other than
$s$.  Intermediate replacement execution is forbidden from entering
$s$ or executing any row whose original target is $s$, since those
rows are all being rewritten.  If every external incoming row has a
nonempty domain, the independently chosen replacements compose and
$s$ can be removed.  Both move directions are considered, so this
also considers repairs which change the predecessor's first head
position by two cells.  It does not cover arbitrary simultaneous
changes to other helper rows.
\end{enumerate}

For an accepted rewrite, relate an original checkpoint to the
configuration with its folded state label and the same head and tape.
Unchanged transitions preserve this relation directly.  A changed
transition is covered by its positive-block domain, and support
exclusion makes the selected replacement trajectory unchanged by the
other edits.  At its endpoint the relation holds again.  This gives
the required finite local simulation proof.  The checker is ordinary
C++ code with this mathematical soundness argument; it is not a
proof-assistant extraction or a kernel-checked formal proof.

\subsection{Recorded bounded results}

All negative runs in Table~\ref{tab:fold-search-results} use
$K=9$, hence $19$ tape cells, and $F=101$.  The tables use inherited
labels, start state $95$, and halt label $254$.  Each has $240$
working states and therefore $964$ candidate edges per row.  An
``accepted row domain'' means one row/edge pair that passes all of
its abstract masks; it need not combine with any other domain to
remove a state.

\begin{table}[htbp]
\centering
\begin{tabular}{llrrr}
Table & Mode & Mask executions & Row domains & Reductions \\
\hline
rotation240 & Pair   & 23,231,947 & 1,030 & 0 \\
rotation240 & Block  & 25,523,988 & 1,101 & 0 \\
rotation240 & Delete & 2,610,950  & 63    & 0 \\
boundary240 & Pair   & 23,227,610 & 1,061 & 0 \\
boundary240 & Block  & 25,578,178 & 1,137 & 0 \\
boundary240 & Delete & 2,671,284  & 68    & 0 \\
\hline
\end{tabular}
\caption{Exact enumeration within the specified local certificate
classes.  Mask executions count attempted edge/mask simulations;
failed edge candidates stop at their first counterexample.}
\label{tab:fold-search-results}
\end{table}

The rotation240 closure contains $11,819,659$ abstract state/window
configurations; the repaired boundary240 closure contains $11,853,163$.
The input SHA-256 digests are:
\begin{quote}\small\ttfamily
rotation240:\par
547dd69da66f66dd48238a53d34b52a3ce85bb6c7d01534dd1cfb815507f248f\par
boundary240:\par
19cd7f7863a681828fd026adbf8ce91c21e3ac7604439bc9d4603ba58661853e
\end{quote}

The positive control uses rotated241 with $K=6$ and $F=31$.
Its closure contains $299,617$ abstract configurations.  Pair mode
finds exactly one reduction, identifying states $72$ and $90$ with
the common rows
\[
  0\mapsto (0,L,144),\qquad 1\mapsto (1,L,36).
\]
This independently checks that the implementation can recognize the
known tape-context fold, rather than merely rejecting all candidates.

\subsection{Reproduction and limits}

The single-file driver \texttt{rerun\_fold\_search.py} embeds its C++17
engine.  It requires Python~3.9 or newer and an available C++17 compiler,
uses no third-party Python packages, and does not import repository
programs.  It compiles its engine in a temporary directory.  For example,
from the directory containing the driver and a copied inherited-label
table:
\begin{verbatim}
python rerun_fold_search.py zfc_rotation240_old.tm \
  --mode all --radius 9 --fuel 101 --output results_rotation240
\end{verbatim}
The output directory must be new.  The driver writes the table and
driver hashes, exact bounds, counts, discovered reductions, accepted
row domains and their support sets, and a compact summary of the
invariant.  The compact invariant file lists fixed offsets and mask
counts; it is not the full set $I_K$.  Rerunning reconstructs the full
closure before checking the domains.  The included result directories
\texttt{portable\_baseline9}, \texttt{portable\_boundary9}, and
\texttt{portable\_control6} record the runs reported above.

The enumeration does not cover all $239$-state tables, all folds,
all choices of invariant, unbounded replacement programs, arbitrary
three-state quotients, or source-level changes.  In particular,
relabeling an endpoint more flexibly can expose recursive loop
identifications outside the literal endpoint comparison used here.
The negative result therefore says that these relatively small,
locally provable rewrites have been exhausted under the stated
window and fuel bounds.  It gives a reusable search component and
a precise boundary for further work, not an impossibility theorem.
